\section{The State of Robot Manipulation Today}
\label{sec:intro:landscape}

Robot manipulation systems now operate in environments that were long considered beyond the reach of automation.
Robots sort and transport goods through warehouses, assemble products on factory lines, assist in surgical and healthcare settings, prepare food, and execute tasks in hazardous and remote environments.
At the same time, the field has expanded beyond rigid industrial automation toward increasingly general-purpose manipulation capabilities: robots are now expected not only to repeat fixed motions reliably, but also to adapt to novel objects, changing environments, and unforeseen task conditions.
The technological progress enabling this expansion has followed two largely separate but equally influential trajectories.

\header{Engineered pipelines at scale}
In industry, manipulation is deployed through modular pipelines: a perception module localizes objects, a planner sequences motions, and a controller executes them.
This architecture underpins robots operating across warehouses and fulfillment centers \cite{CensusSurveyRobotic2022}, factory assembly and welding lines \cite{factory_efficiency_2020}, surgical and assistive healthcare settings \cite{Kyrarini2021}, chemical and pharmaceutical processing \cite{chem1,chem2}, food service and meal preparation \cite{meal1,assistfeed1,assistfeed2}, and space and extreme-environment operations \cite{papadopoulos2021robotic,gaines2016productivity}.
The defining properties of these systems are reliability and debuggability: each module is maintained by specialists who understand its failure modes, and performance improves through engineering effort accumulated over many deployment cycles.
The result is highly capable systems within their operational envelopes---systems trusted with production-scale throughput and human safety.

\header{General-purpose learned policies}
In parallel, a different paradigm has emerged in academia and is now entering industry.
General-purpose robot policies trained on large, diverse demonstration datasets have demonstrated impressive generalization: a single policy can handle novel objects, environmental variation, and modest task variation without per-task redesign \cite{chi2023diffusionpolicy,team2024octo,o2024open,zitkovich2023rt}.
Reinforcement learning has pushed further still, enabling dexterous in-hand manipulation of unprecedented complexity by training entirely in simulation \cite{andrychowicz2018learning}.
The aspiration of this line of work is human-level dexterity accessible through learned behavior---reducing the per-task engineering burden and enabling continuous improvement as deployment data accumulates.

\header{Shared limitations}
Despite their apparent differences, both paradigms share a set of fundamental limitations.
Both require substantial engineering effort to robustify against real-world variability: pipeline systems need task-specific tuning and edge-case handling at every component; learned systems require careful dataset curation, sim-to-real transfer, and distribution coverage to avoid brittle failures \cite{billard2019trends}.
Both exhibit brittleness at component interfaces---errors between perception and planning, and between planning and control, propagate in ways that are difficult to anticipate or recover from \cite{abderezaei2024clutter,briscoe2024exploring}.
% Both rely on simplified models of physics: classical pipelines use analytical models with bounded validity, while learned policies treat dynamics as an implicit property of demonstrated trajectories rather than as an explicit constraint \cite{yin2021modeling}.
In both paradigms, dynamics are typically treated as secondary: classical pipelines approximate them through simplified analytical models, while learned policies absorb them implicitly through demonstrated behavior rather than representing them as explicit constraints \cite{yin2021modeling}.
And neither generalizes readily to task contexts substantially different from those encountered during development or training \cite{o2024open}.

\header{The geometric assumption}
Underlying both paradigms is a shared assumption that has shaped manipulation research for decades: that manipulation is fundamentally a geometric problem.
The dominant objects of reasoning are collision clearance, grasp pose, reachability, and singularity avoidance---broadly, kinematic constraints on robot configuration.
This assumption is not arbitrary.
Grasping has attracted disproportionate research attention \cite{billard2019trends,bohg2013data,mahler2017dex} precisely because form-closure reduces object state to a geometry problem: once an object is secured, its state is predictable and planning reduces to obstacle avoidance \cite{mason1986mechanics,bullock2011classifying}.
Recent years have brought dramatic improvements to geometric planning infrastructure---GPU-parallelized collision checking and forward kinematics \cite{sundaralingam2023curobo,thomason2024motions}, large-scale motion datasets \cite{o2024open}, and high-fidelity simulation environments \cite{coumans2021}---but the underlying planning abstraction has remained largely geometric.

\header{The dynamic gap}
The cost of this assumption becomes visible at the boundary of kinematic behaviors.
A broad class of manipulation tasks is not adequately captured by geometry alone: tasks where joint torques, object inertia, contact impulses, and fluid dynamics jointly determine the outcome.
Non-prehensile rearrangement through impulsive contact \cite{pasricha2022pokerrt,mason2018toward}, transport of liquid-filled containers \cite{abderezaei2024clutter}, high-speed inertial transport \cite{ichnowski2022gomp,zeng2020tossingbot}, operation near or beyond rated payload capacity, and adaptation to mid-task mechanical failures \cite{briscoe2024exploring} all fall into this category.
These tasks involve two distinct kinds of dynamic constraint: \textit{higher-order robot constraints}---bounds on velocity, acceleration, jerk, and torque that arise from the robot's own kinematics and actuator limits \cite{pham2018new,berscheid2021jerk}---and \textit{task-induced dynamic constraints} that arise from the physics of robot-object interaction, including contact mechanics, inertial coupling, and motion-induced phenomena such as sloshing.
Non-prehensile manipulation, which encompasses this class of behaviors, has received comparatively little attention relative to grasping \cite{ruggiero2018nonprehensile,mason2018toward}, and the gap between what current systems can do and what these tasks demand is consequential.

% \header{The argument}
% The central claim of this thesis is that accessing the full range of manipulation capability afforded by modern robot embodiments---beyond the kinematic behaviors that dominate current systems---requires explicitly modeling and integrating dynamic constraints into the motion generation process.
% Tasks involving complex object interactions, unstructured environments, and high-speed or high-load operations cannot be adequately planned or controlled using geometry alone.
% Dynamics must be treated not as implicit background or a post-hoc check, but as structured information to be modeled, learned, and exploited.

% \header{A spectrum of approaches}
% How to incorporate dynamics into manipulation planning is not a settled question, and the literature reflects a spectrum of strategies.
% At one end, \textit{plan-and-filter} methods generate geometric paths and apply time parameterization or post-processing to enforce dynamic feasibility, iterating until constraints are satisfied \cite{pham2018new,berscheid2021jerk,kuntz2019fast}.
% \textit{Kinodynamic search} methods embed dynamics directly in the planning search, propagating controls through physics-aware forward models to build dynamically feasible trees \cite{pasricha2022pokerrt,li2016asymptotically,nayak2022bidirectional}.
% \textit{Constrained trajectory optimization} formulates dynamics as explicit equality or inequality constraints and solves for a feasible trajectory from a given seed \cite{schulman2013finding,kalakrishnan2011stomp,ichnowski2020gomp}.
% \textit{Learned motion generation} approaches---including diffusion-based trajectory models and foundation policies---implicitly encode dynamics through training data and conditioning signals, avoiding runtime constraint evaluation altogether \cite{chi2023diffusionpolicy,carvalho2023motion,saha2024edmp,fishman2023motion,qureshi2020motion}.
% Finally, \textit{hybrid} approaches combine these paradigms, for example using learned initializations to seed optimization or using sampling to guide learning \cite{kuntz2019fast,ichnowski2020deep,huang2024diffusionseeder}.
% Each strategy makes different tradeoffs along four recurring dimensions: modeling burden, problem dimensionality, system integration complexity, and practical usability---the axes that organize the prior work reviewed next.
